%%%%%%%%%%%%%%%%%%%%%%%%%%%%%%%%%%%%%%%%%%%%%%%%%%%%%%%%%%%%%%%%%%%%%%
% Lớp 9 | Chương V. Đường tròn | Vị trí tương đối của hai đường tròn
% Dạng 9H5?17-2: Tiếp tuyến chung, bài toán tổng hợp
%%%%%%%%%%%%%%%%%%%%%%%%%%%%%%%%%%%%%%%%%%%%%%%%%%%%%%%%%%%%%%%%%%%%%%


% nguồn: 2026-11-21_lop9_chuong5_bai14_17 (nhập 21/11/2026)
\begin{ex}%[9H5N17-2]%[ma:eac82194]
% nguon: SGK Toán 9 KNTT tập 1 – Bài 17, Vị trí tương đối của hai đường tròn
	Hai đường tròn ở ngoài nhau có tất cả bao nhiêu tiếp tuyến chung?
	\choice
	{$1$}
	{$2$}
	{$3$}
	{\True $4$}
	\loigiai{
		Hai đường tròn ở ngoài nhau có $2$ tiếp tuyến chung ngoài và $2$ tiếp tuyến chung trong, tổng cộng là $4$ tiếp tuyến chung.
	}
\end{ex}

% nguồn: 2026-11-21_lop9_chuong5_bai14_17 (nhập 21/11/2026)
\begin{ex}%[9H5H17-2]%[ma:d5f482cb]
% nguon: SGK Toán 9 KNTT tập 1 – Bài 17, Vị trí tương đối của hai đường tròn
	Cho hai đường tròn $(O; 9\text{ cm})$ và $(O'; 4\text{ cm})$ tiếp xúc ngoài tại điểm $A$. Gọi $BC$ là một tiếp tuyến chung ngoài của hai đường tròn ($B \in (O)$, $C \in (O')$). Xét tính đúng, sai của các khẳng định sau:
	\choiceTF
	{\True Độ dài đoạn nối tâm $OO'$ bằng $13\text{ cm}$}
	{\True Tứ giác $OBCO'$ là hình thang vuông}
	{Độ dài đoạn tiếp tuyến chung ngoài $BC$ bằng $10\text{ cm}$}
	{\True Diện tích tứ giác $OBCO'$ bằng $78\text{ cm}^2$}
	\loigiai{
		\begin{itemchoice}
			\itemch Đúng. Vì hai đường tròn tiếp xúc ngoài nên $OO' = R + r = 9 + 4 = 13\text{ cm}$.
			\itemch Đúng. Vì $OB \perp BC$ và $O'C \perp BC$ nên $OB \parallel O'C$ và $\widehat{OBC} = \widehat{O'CB} = 90^\circ$, do đó $OBCO'$ là hình thang vuông.
			\itemch Sai. Kẻ $O'H \perp OB$ tại $H$, khi đó $HBCO'$ là hình chữ nhật nên $BC = O'H$ và $OH = OB - O'C = 9 - 4 = 5\text{ cm}$. Áp dụng định lí Pythagore cho tam giác $OHO'$ vuông tại $H$: $BC = O'H = \sqrt{OO'^2 - OH^2} = \sqrt{13^2 - 5^2} = 12\text{ cm}$.
			\itemch Đúng. Diện tích hình thang vuông $OBCO'$ bằng $\dfrac{OB + O'C}{2} \cdot BC = \dfrac{9 + 4}{2} \cdot 12 = 78\text{ cm}^2$.
		\end{itemchoice}
	}
\end{ex}

% nguồn: 2026-11-21_lop9_chuong5_bai14_17 (nhập 21/11/2026)
\begin{bt}%[9H5V17-2]%[ma:46551519]
% nguon: SGK Toán 9 KNTT tập 1 – Bài 17, Vị trí tương đối của hai đường tròn
	Hai bánh răng hình tròn có bán kính $R = 8\text{ cm}$ và $r = 2\text{ cm}$ đặt tiếp xúc ngoài nhau. Một đoạn dây truyền động căng thẳng tiếp xúc ngoài với cả hai bánh răng tại hai điểm $A$ và $B$. Tính độ dài đoạn thẳng $AB$ (đơn vị cm).
	\shortans{8}
	\loigiai{
		Gọi $O$ và $O'$ lần lượt là tâm hai bánh răng bán kính $R = 8\text{ cm}$ và $r = 2\text{ cm}$.\\
		Vì hai đường tròn tiếp xúc ngoài nên $OO' = R + r = 8 + 2 = 10\text{ cm}$.\\
		Kẻ $O'H \perp OA$ tại $H$, khi đó $HABO'$ là hình chữ nhật nên $AB = O'H$ và $OH = R - r = 8 - 2 = 6\text{ cm}$.\\
		Trong tam giác $OHO'$ vuông tại $H$: $AB = O'H = \sqrt{OO'^2 - OH^2} = \sqrt{10^2 - 6^2} = 8\text{ cm}$.
	}
\end{bt}

% nguồn: 2026-11-21_lop9_chuong5_bai16_17_bs (nhập 21/11/2026)
\begin{ex}%[9H5N17-2]%[ma:d3a8f652]
% nguon: SGK Toán 9 KNTT tập 1 – Bài 17, Vị trí tương đối của hai đường tròn
	Hai đường tròn tiếp xúc ngoài có tất cả bao nhiêu tiếp tuyến chung?
	\choice
	{$1$}
	{$2$}
	{\True $3$}
	{$4$}
	\loigiai{
		Hai đường tròn tiếp xúc ngoài có $2$ tiếp tuyến chung ngoài và $1$ tiếp tuyến chung trong đi qua tiếp điểm, tổng cộng là $3$ tiếp tuyến chung.
	}
\end{ex}

% nguồn: 2026-11-21_lop9_chuong5_bai16_17_bs (nhập 21/11/2026)
\begin{ex}%[9H5N17-2]%[ma:41272858]
% nguon: SGK Toán 9 KNTT tập 1 – Bài 17, Vị trí tương đối của hai đường tròn
	Hai đường tròn cắt nhau có tất cả bao nhiêu tiếp tuyến chung?
	\choice
	{$0$}
	{$1$}
	{\True $2$}
	{$3$}
	\loigiai{
		Hai đường tròn cắt nhau chỉ có $2$ tiếp tuyến chung ngoài, không có tiếp tuyến chung trong.
	}
\end{ex}

% nguồn: 2026-11-21_lop9_chuong5_bai16_17_bs (nhập 21/11/2026)
\begin{ex}%[9H5H17-2]%[ma:671096f1]
% nguon: SGK Toán 9 KNTT tập 1 – Bài 17, Vị trí tương đối của hai đường tròn
	Cho hai đường tròn $(O; 16\text{ cm})$ và $(O'; 9\text{ cm})$ tiếp xúc ngoài tại điểm $A$. Gọi $DE$ là một tiếp tuyến chung ngoài của hai đường tròn ($D \in (O)$, $E \in (O')$). Xét tính đúng, sai của các khẳng định sau:
	\choiceTF
	{\True Độ dài đoạn nối tâm $OO'$ bằng $25\text{ cm}$}
	{\True Đường thẳng đi qua $A$ và vuông góc với $OO'$ là tiếp tuyến chung trong của hai đường tròn}
	{Độ dài đoạn tiếp tuyến chung ngoài $DE$ bằng $20\text{ cm}$}
	{\True Diện tích hình thang vuông $ODEO'$ bằng $300\text{ cm}^2$}
	\loigiai{
		\begin{itemchoice}
			\itemch Đúng. Vì hai đường tròn tiếp xúc ngoài nên $OO' = 16 + 9 = 25\text{ cm}$.
			\itemch Đúng. Đường thẳng qua $A$ vuông góc với $OO'$ tại $A$ vuông góc với cả bán kính $OA$ và $O'A$ nên là tiếp tuyến chung trong của hai đường tròn.
			\itemch Sai. Kẻ $O'H \perp OD$ tại $H$, ta có $OH = 16 - 9 = 7\text{ cm}$ nên $DE = O'H = \sqrt{25^2 - 7^2} = \sqrt{576} = 24\text{ cm}$.
			\itemch Đúng. Diện tích hình thang vuông $ODEO'$ là $\dfrac{16 + 9}{2} \cdot 24 = 300\text{ cm}^2$.
		\end{itemchoice}
	}
\end{ex}

% nguồn: 2026-11-21_lop9_chuong5_bai16_17_bs (nhập 21/11/2026)
\begin{ex}%[9H5H17-2]%[ma:3c1604cb]
% nguon: SGK Toán 9 KNTT tập 1 – Bài 17, Vị trí tương đối của hai đường tròn
	Cho hai đường tròn $(O; 7\text{ cm})$ và $(O'; 3\text{ cm})$ có khoảng cách hai tâm $OO' = 15\text{ cm}$. Xét tính đúng, sai của các khẳng định sau:
	\choiceTF
	{\True Hai đường tròn $(O)$ và $(O')$ ở ngoài nhau}
	{\True Hai đường tròn có $2$ tiếp tuyến chung ngoài và $2$ tiếp tuyến chung trong}
	{Khoảng cách ngắn nhất giữa một điểm trên $(O)$ và một điểm trên $(O')$ bằng $4\text{ cm}$}
	{\True Độ dài đoạn tiếp tuyến chung trong nối hai tiếp điểm của hai đường tròn bằng $5\sqrt{5}\text{ cm}$}
	\loigiai{
		\begin{itemchoice}
			\itemch Đúng. Vì $OO' = 15 > 7 + 3 = 10$ nên hai đường tròn ở ngoài nhau.
			\itemch Đúng. Hai đường tròn ở ngoài nhau luôn có $2$ tiếp tuyến chung ngoài và $2$ tiếp tuyến chung trong.
			\itemch Sai. Khoảng cách ngắn nhất giữa hai điểm trên hai đường tròn là $OO' - R - r = 15 - 7 - 3 = 5\text{ cm}$.
			\itemch Đúng. Độ dài đoạn tiếp tuyến chung trong là $\sqrt{OO'^2 - (R + r)^2} = \sqrt{15^2 - 10^2} = \sqrt{125} = 5\sqrt{5}\text{ cm}$.
		\end{itemchoice}
	}
\end{ex}

% nguồn: 2026-11-21_lop9_chuong5_bai16_17_bs (nhập 21/11/2026)
\begin{bt}%[9H5V17-2]%[ma:7362b47e]
% nguon: SGK Toán 9 KNTT tập 1 – Bài 17, Vị trí tương đối của hai đường tròn
	Cho hai đường tròn $(O; 10\text{ cm})$ và $(O'; 2\text{ cm})$ ở ngoài nhau có khoảng cách hai tâm $OO' = 17\text{ cm}$. Tính độ dài đoạn tiếp tuyến chung ngoài nối hai tiếp điểm của hai đường tròn đó (đơn vị cm).
	\shortans{15}
	\loigiai{
		Gọi $A \in (O)$ và $B \in (O')$ là hai tiếp điểm của tiếp tuyến chung ngoài. Kẻ $O'H \perp OA$ tại $H$, khi đó $HABO'$ là hình chữ nhật nên $AB = O'H$ và $OH = R - r = 10 - 2 = 8\text{ cm}$.\\
		Trong tam giác $OHO'$ vuông tại $H$: $AB = O'H = \sqrt{OO'^2 - OH^2} = \sqrt{17^2 - 8^2} = \sqrt{225} = 15\text{ cm}$.
	}
\end{bt}

% nguồn: 2026-11-21_lop9_chuong5_bai16_17_bs (nhập 21/11/2026)
\begin{bt}%[9H5V17-2]%[ma:85d72255]
% nguon: SGK Toán 9 KNTT tập 1 – Bài 17, Vị trí tương đối của hai đường tròn
	Cho hai đường tròn $(O; 6\text{ cm})$ và $(O'; 3\text{ cm})$ ở ngoài nhau có khoảng cách hai tâm $OO' = 15\text{ cm}$. Tính độ dài đoạn tiếp tuyến chung trong nối hai tiếp điểm của hai đường tròn đó (đơn vị cm).
	\shortans{12}
	\loigiai{
		Gọi $A \in (O)$ và $B \in (O')$ là hai tiếp điểm của tiếp tuyến chung trong. Kẻ $O'K \perp OA$ tại $K$ (trên tia đối của tia $AO$), ta có $OK = R + r = 6 + 3 = 9\text{ cm}$ và $AB = O'K$.\\
		Trong tam giác $OKO'$ vuông tại $K$: $AB = O'K = \sqrt{OO'^2 - OK^2} = \sqrt{15^2 - 9^2} = \sqrt{144} = 12\text{ cm}$.
	}
\end{bt}
